\documentclass[10pt]{article}

% ── Core packages ─────────────────────────────────────────────────────────────
\usepackage[utf8]{inputenc}
\usepackage[T1]{fontenc}
\usepackage{lmodern}
\usepackage{microtype}

% ── Page geometry (NeurIPS/ICML style: 6.5in × 9in text block) ───────────────
\usepackage[
  top=1in, bottom=1in,
  left=0.75in, right=0.75in,
  columnsep=0.25in
]{geometry}
\setlength{\parindent}{0pt}
\setlength{\parskip}{4pt}

% ── Mathematics ───────────────────────────────────────────────────────────────
\usepackage{amsmath,amssymb,amsthm}

% ── Figures and tables ────────────────────────────────────────────────────────
\usepackage{graphicx}
\usepackage{booktabs}
\usepackage{multirow}
\usepackage{float}
\usepackage{caption}
\usepackage{subcaption}

% ── Hyperlinks and references ─────────────────────────────────────────────────
\usepackage[colorlinks=true,linkcolor=blue,citecolor=blue,urlcolor=blue]{hyperref}
\usepackage{natbib}

% ── Utilities ─────────────────────────────────────────────────────────────────
\usepackage{xcolor}
\usepackage{url}
\usepackage{enumitem}

% ─────────────────────────────────────────────────────────────────────────────

\begin{document}

\title{Do Published Attention-Block Gains Survive Common Image Corruptions?\\
A Matched-Budget CIFAR-10 / CIFAR-10-C Stress Test of SE, BAM, and CBAM}
\author{Anonymous Submission}
\date{}
\maketitle

\begin{abstract}
Channel and spatial attention blocks are compared almost exclusively on clean
test accuracy, and one of them positions itself as an improvement over another
on precisely that basis. Whether the ordering this induces survives the
distribution shift a deployed classifier meets has not been tested for this
method family under a matched protocol. We train squeeze-and-excitation, the
bottleneck attention module, and the convolutional block attention module, plus
an attention-free baseline, into one shared CIFAR-scale residual backbone at an
identical budget, inserting each module at the identical point so that only the
module varies, and evaluate the resulting twelve checkpoints on CIFAR-10-C
without retraining. The clean-accuracy ordering
$\text{CBAM} > \text{SE} > \text{none} > \text{BAM}$ is preserved at every
corruption severity, on mean corruption accuracy, and on a
starting-point-normalised relative degradation measure: for this method set the
field's single metric was not misleading. Two secondary findings complicate the
picture. Clean accuracy cannot separate the top two arms ($0.11$ pp against a
combined across-seed spread of $0.26$ pp) whereas mean corruption accuracy can
($0.92$ pp against $0.49$ pp), so the corruption axis is here the more
discriminative measurement rather than the noisier one; and the parallel
channel-plus-spatial module falls $4.5$ pp \emph{below} the attention-free
baseline. We also report a discrepancy between that module's published fusion
equation and the fusion its official code computes.
\end{abstract}

% ── 1. Introduction ───────────────────────────────────────────────────────────
\section{Introduction}
\label{sec:introduction}

Attention blocks are among the cheapest architectural additions in modern
computer vision. A handful of extra parameters inserted into an existing
backbone is reported to buy a consistent accuracy gain, and
squeeze-and-excitation (SE)~\citep{hu2018squeeze}, the bottleneck attention
module (BAM)~\citep{park2018bam}, and the convolutional block attention module
(CBAM)~\citep{woo2018cbam} have become canonical instances of the
pattern~\citep{guo2022attention}. Their reported gains are the reason the
applied literature routinely bolts them onto convolutional networks.

Those gains are all reported on clean test data. Each of these methods is
evaluated on top-1 and top-5 accuracy, and CBAM positions itself explicitly as
an improvement over SE's channel-only design on that basis. None reports how
the resulting models behave when the input distribution shifts.

This matters because the ordering these papers induce is what a practitioner
actually consumes. Someone selecting an attention block for a deployed
classifier reads off a ranking measured under conditions the deployment will
not reproduce: sensors defocus, illumination drifts, contrast collapses, images
are recompressed in transit. Convolutional classifiers are known to be brittle
to exactly these
perturbations~\citep{hendrycks2019benchmarking,geirhos2019imagenet}, and a
substantial literature now treats corruption robustness as a first-class
evaluation axis~\citep{hendrycks2020augmix}. The question that follows is
therefore not whether attention blocks improve clean accuracy, but whether the
\emph{ordering} they induce on clean data is the ordering that survives
corruption.

The existing literature cannot answer this, for two reasons. First, each method
is validated in its own backbone, at its own scale, under its own training
recipe, so a cross-paper accuracy difference conflates the module with the
architecture and the optimisation schedule. Second, and more subtly, BAM and
CBAM originate from a single shared public implementation and inherit a common
insertion convention that an independently reimplemented SE does not; even a
same-codebase comparison therefore confounds module identity with module
placement.

We remove both confounds. All three attention blocks, plus an attention-free
baseline, are trained into one shared CIFAR-scale residual
backbone~\citep{he2016deep}, with each module inserted at the identical point in
the identical block, under an identical optimiser, schedule, augmentation
policy, epoch budget and seed set. Because the module is the only term that
varies, any measured difference is attributable to the module itself. We then
evaluate the trained checkpoints, without retraining, on a common-corruption
benchmark that shares the clean dataset's exact label space and
resolution~\citep{hendrycks2019benchmarking,krizhevsky2009learning}, and re-rank
the arms by mean corruption accuracy and by a relative degradation measure
rather than by clean top-1.

Our hypothesis was that the ranking might not be invariant. Channel-only
recalibration derives its gate from a globally pooled descriptor, whereas the
spatial branches of BAM and CBAM compute attention from local feature
statistics. Spatial masks depend on where informative structure lies in the
image, and corruptions such as defocus blur and contrast reduction attack that
structure directly, so a spatial gate that mislocalises under corruption could
degrade faster than a channel-only gate. We stressed in advance that the
competing outcome is equally informative, and committed to reporting it as
prominently. That is the outcome we obtained.

Our contributions are:

\begin{itemize}[leftmargin=1.5em]
  \item \textbf{A matched-budget, single-backbone comparison} of three published
    attention blocks against an attention-free baseline, in which backbone,
    insertion point, optimiser, schedule, augmentation and seed set are all held
    fixed, so that differences are attributable to the module rather than to its
    host or its placement.
  \item \textbf{A negative result, reported as the primary finding.} The
    clean-accuracy ordering is preserved at every corruption severity, on mean
    corruption accuracy, and on relative degradation. For this method set,
    ranking on clean top-1 was not misleading.
  \item \textbf{A measurement-power observation that qualifies it.} Clean
    accuracy cannot resolve the top two arms, while mean corruption accuracy
    can. The corruption axis is here the more discriminative measurement, which
    is a reason to report it even when it does not reorder anything.
  \item \textbf{A provenance record}, including a discrepancy between BAM's
    published fusion equation and the fusion its official public code computes,
    which to our knowledge has not previously been reported.
\end{itemize}

Section~\ref{sec:related} positions the work. Section~\ref{sec:method} describes
the backbone, the modules and their provenance, and the fidelity checks.
Section~\ref{sec:setup} gives the protocol and metrics. Section~\ref{sec:results}
reports results and answers the title question. Section~\ref{sec:limitations}
scopes the claim and Section~\ref{sec:conclusion} concludes.

% ── 2. Related Work ───────────────────────────────────────────────────────────
\section{Related Work}
\label{sec:related}

\subsection{Channel and Spatial Attention Blocks}

Squeeze-and-excitation~\citep{hu2018squeeze} introduced channel-only feature
recalibration: global average pooling reduces a feature map to a per-channel
descriptor, a two-layer bottleneck multilayer perceptron maps that descriptor to
per-channel gates, a sigmoid bounds them, and the features are rescaled
channel-wise. The design is cheap, architecture-agnostic, and widely adopted;
subsequent work has refined the channel-gating mechanism itself, for example by
replacing the bottleneck with a local cross-channel
interaction~\citep{wang2020eca}.

BAM~\citep{park2018bam} and CBAM~\citep{woo2018cbam} extend this by adding a
spatial branch. BAM computes channel and spatial attention in parallel and fuses
them; CBAM applies channel attention and then spatial attention in sequence, and
pools both average and maximum statistics in its channel gate. Both spatial
branches derive their attention maps from local feature statistics rather than
from a global descriptor, which is the property our hypothesis turned on. Broad
surveys confirm the pattern across the family~\citep{guo2022attention}.

What unites this cluster for our purposes is its evaluation protocol. The
metrics reported are clean top-1 and top-5 classification accuracy, with mean
average precision added where detection is evaluated. No member of this cluster
reports a robustness or distribution-shift metric, so the relative ordering
these papers establish has never been tested away from the clean distribution.

\subsection{Corruption Robustness Benchmarking}

A parallel literature evaluates classifiers under synthetic but realistic
corruptions applied at graded severities~\citep{hendrycks2019benchmarking}. The
protocol fixes a set of corruption types, applies each at five severity levels,
and summarises degradation across the grid, which makes it possible to compare
how fast models degrade rather than only where they start. Related work
established that convolutional classifiers rely on cues these corruptions
disrupt~\citep{geirhos2019imagenet}, and that targeted augmentation can mitigate
the damage~\citep{hendrycks2020augmix}.

The corrupted benchmark used here shares the clean dataset's exact label space
and resolution~\citep{krizhevsky2009learning}, so evaluating a trained model on
it requires inference only: no retraining, no relabelling, no architectural
change. That is what makes the audit cheap enough to be worth running, and makes
its absence from the attention-block literature harder to explain.

Characteristically, however, this literature deploys corruption benchmarks to
\emph{validate a newly proposed intervention}. The benchmark is the evidence
that a new method helps, rather than the instrument for auditing the ranking of
methods that already exist.

\subsection{Attention Mechanisms Under Distribution Shift}

Recent work sits at the intersection. Attention-guided repair procedures improve
the robustness of convolutional networks against common
corruptions~\citep{ar2repair2025}, and channel attention has been combined with
multi-scale push--pull inhibition to the same end~\citep{pushpull2025}. Both are
validated on corruption benchmarks of the kind we use.

Both, however, are proposals: each introduces a new mechanism and uses the
corruption benchmark as evidence that it works. Neither is a head-to-head
comparison of the already-published blocks under a single matched training
protocol with a fixed insertion point. That controlled comparison, and the
ranking-stability question it makes answerable, is the gap this paper fills.

% ── 3. Method ─────────────────────────────────────────────────────────────────
\section{Method}
\label{sec:method}

\subsection{Shared Backbone and a Single Insertion Point}
\label{sec:backbone}

All arms use one CIFAR-adapted ResNet-18: a $3\times3$ stride-1 stem with no
initial max-pool, preserving $32\times32$ spatial detail, followed by four
stages of two basic blocks at widths $64/128/256/512$, global average pooling,
and a linear classifier. Within each basic block the attention module is applied
to the residual branch after the second batch
normalisation~\citep{ioffe2015batch} and before the residual addition:
%
\begin{equation}
  \mathrm{out} = \mathrm{ReLU}\!\left(
    A\big(\mathrm{BN}_2(\mathrm{conv}_2(\mathrm{ReLU}(\mathrm{BN}_1(\mathrm{conv}_1(x)))))\big)
    + \mathrm{shortcut}(x)\right),
  \label{eq:block}
\end{equation}
%
where $A$ is the attention module under test, or the identity for the baseline
arm. The identity shortcut never passes through $A$. Figure~\ref{fig:backbone}
shows the block.

Because $A$ is the only term in Equation~\ref{eq:block} that varies, any
measured difference is attributable to the module rather than to placement, host
architecture, or training recipe. This is the study's primary confound control.
It has a cost, which we state plainly: no arm reproduces its own paper's
placement --- official BAM sits \emph{between} stages rather than inside blocks
--- so absolute accuracies here are not comparable to published numbers, and we
claim none. Only the within-study ordering is claimed.

The arms differ only by attention parameters: $11.17$M for the baseline,
$11.26$M for SE, $11.36$M for BAM, and $11.26$M for CBAM.

\begin{figure}[t]
  \centering
  \includegraphics[width=0.74\linewidth]{figures/fig_shared_backbone_overview.png}
  \caption{The shared residual block and its single attention slot. Every arm
  uses this identical block; only the module $A$ changes, and the identity
  shortcut never passes through it. Holding the insertion point fixed is what
  makes a measured difference attributable to the module rather than to its
  placement.}
  \label{fig:backbone}
\end{figure}

\subsection{Attention Modules and Their Provenance}
\label{sec:modules}

\textbf{SE} is reimplemented from Equations 2--3 of~\citet{hu2018squeeze}:
global average pool to $z \in \mathbb{R}^{C}$, then
$s = \sigma(W_2\,\delta(W_1 z))$ with $W_1 \in \mathbb{R}^{(C/r) \times C}$,
then channel-wise rescale by $s$. This is the study's one reimplementation, and
it is necessary: the official release is Caffe-only, consisting of
\texttt{.prototxt} model definitions with custom \texttt{.cpp}/\texttt{.cu}
layers, and contains no Python. The block is fully specified by two equations,
so the reimplementation is unambiguous. We use $r=16$, the paper's default, and
record it as one point in the design space rather than a re-derivation of a
tuned value.

\textbf{BAM} and \textbf{CBAM} are ported directly from the module definitions
of the official public implementation. We take only the module definitions; the
repository's own ImageNet-scale model file is not used. Layer types, ordering,
pooling choices, bias settings, batch-norm placement and hyperparameters,
$C/r$ bottleneck widths, channel-pool concatenation order and BAM's dilation-4
convolution stack are all preserved. The only change is replacing a deprecated
sigmoid entry point removed from modern versions of the framework; the
repository has had no commits since March 2023 and ships no dependency pins.

Two details are load-bearing, because conflating them would silently weaken the
comparison.

\emph{BAM's and CBAM's channel gates are different mechanisms.} Both files
define a class of the same name, but BAM's pools average statistics only and
places a one-dimensional batch normalisation between its MLP layers, while
CBAM's pools average \emph{and} maximum statistics through a shared MLP with no
normalisation. Sharing one implementation between the two arms would partially
collapse the channel-only-versus-channel-plus-spatial contrast this study
measures.

\emph{BAM's paper and its official code disagree on the fusion operator.}
\citet{park2018bam} specify additive fusion,
$M(F) = \sigma(M_c(F) + M_s(F))$. The official code instead computes
$\mathrm{att} = 1 + \sigma(M_c(F) \cdot M_s(F))$ and returns
$\mathrm{att} \cdot F$: a multiplicative fusion with a residual offset. Since
the artefact this literature actually reuses is the code rather than the
equation, we follow the code, and we report the discrepancy rather than
resolving it silently. To our knowledge it has not previously been noted.

\subsection{Mechanical Verification of Module Fidelity}
\label{sec:verification}

The comparison's validity rests entirely on the modules being what they claim to
be, so fidelity is asserted mechanically rather than by inspection. A
verification suite of 54 checks asserts structural properties --- BAM's
average-only pooling and one-dimensional batch normalisation, CBAM's dual
pooling and batch-normalised $7\times7$ spatial convolution, SE's pooling,
bottleneck width and sigmoid gate --- and, critically, behavioural ones. BAM's
output-to-input ratio must lie in $(1,2)$, as $1 + \sigma(\cdot)$ requires, and
CBAM's in $(0,1)$, as a pure multiplicative gate requires. Fusion semantics are
thus verified by behaviour rather than by reading source code. The suite also
asserts that the two channel-gate classes are distinct objects, so the defect
described in Section~\ref{sec:modules} cannot silently return. It runs before
every training installment.

% ── 4. Experimental Setup ─────────────────────────────────────────────────────
\section{Experimental Setup}
\label{sec:setup}

\subsection{Data and Training Protocol}

Training and clean evaluation use CIFAR-10~\citep{krizhevsky2009learning}:
$50{,}000$ training and $10{,}000$ test images at $32\times32$ across 10
classes. Corrupted evaluation uses CIFAR-10-C~\citep{hendrycks2019benchmarking},
restricted to four corruption types --- brightness, contrast, defocus blur and
elastic transform --- at severities 1 through 5, with $10{,}000$ images per type
and severity. Test-time preprocessing is identical for clean and corrupted
images.

Every arm receives an identical recipe: SGD with learning rate $0.1$, momentum
$0.9$, weight decay $5\times10^{-4}$, batch size $128$, cosine
annealing~\citep{loshchilov2017sgdr}, and standard augmentation (random
$32\times32$ crop with 4-pixel padding, random horizontal flip). The budget is
10 epochs per run and the seed set is $\{1,2,3\}$, giving 12 training runs. No
arm uses a pretrained or externally obtained checkpoint. Mixed
precision~\citep{micikevicius2018mixed} is used for training; validation runs in
full precision so that checkpoint selection and final evaluation share one
numeric path. All runs used a single NVIDIA RTX A2000 with 6\,GB of memory.

\subsection{Metrics}

We report three. \emph{Clean top-1 accuracy} is the metric the field currently
ranks on and the severity-zero reference point. \emph{Mean corruption accuracy}
(mCA) at each severity averages over corruption types. \emph{Relative robustness
drop}, $(\text{clean} - \text{corrupted})/\text{clean}$, normalises out each
arm's clean starting point so that an arm beginning lower is not mistaken for a
robust one merely because it had less accuracy to lose.

All results are averaged over seeds and reported with the across-seed standard
deviation. Before any ordering is claimed, we check whether the gap between the
top two arms exceeds their combined across-seed spread. We compare against the
spread of the two arms being compared, not the largest spread anywhere in the
grid; the latter belongs to the higher-variance baseline and BAM arms and would
wrongly declare a resolvable difference unresolved.

% ── 5. Results ────────────────────────────────────────────────────────────────
\section{Results and Analysis}
\label{sec:results}

\subsection{Accuracy Across Corruption Severity}

Table~\ref{tab:main} reports every arm at every severity, with mCA and relative
drop. Figure~\ref{fig:severity} plots the same data.

\begin{table}[t]
  \centering
  \caption{Top-1 accuracy (\%) by corruption severity, averaged over four
  corruption types and three seeds, with across-seed standard deviation. mCA is
  the mean over severities 1--5; relative drop is
  $(\text{clean}-\text{mCA})/\text{clean}$, where lower is better.}
  \label{tab:main}
  \small
  \begin{tabular}{lcccccccc}
    \toprule
    Model & Clean & Sev.\ 1 & Sev.\ 2 & Sev.\ 3 & Sev.\ 4 & Sev.\ 5 & mCA & Rel.\ drop \\
    \midrule
    none & $86.09${\scriptsize$\pm2.10$} & $83.96${\scriptsize$\pm2.25$} & $80.99${\scriptsize$\pm2.93$} & $75.53${\scriptsize$\pm3.00$} & $67.12${\scriptsize$\pm2.84$} & $55.01${\scriptsize$\pm0.99$} & $72.52${\scriptsize$\pm2.31$} & $15.77\%$ \\
    SE   & $89.21${\scriptsize$\pm0.18$} & $87.29${\scriptsize$\pm0.10$} & $85.05${\scriptsize$\pm0.05$} & $80.09${\scriptsize$\pm0.23$} & $71.98${\scriptsize$\pm0.06$} & $57.16${\scriptsize$\pm0.05$} & $76.31${\scriptsize$\pm0.04$} & $14.45\%$ \\
    BAM  & $81.56${\scriptsize$\pm2.47$} & $79.28${\scriptsize$\pm2.71$} & $75.51${\scriptsize$\pm3.09$} & $70.08${\scriptsize$\pm3.08$} & $62.83${\scriptsize$\pm2.67$} & $52.91${\scriptsize$\pm2.20$} & $68.12${\scriptsize$\pm2.54$} & $16.49\%$ \\
    CBAM & $\mathbf{89.32}${\scriptsize$\pm0.08$} & $\mathbf{87.68}${\scriptsize$\pm0.22$} & $\mathbf{85.58}${\scriptsize$\pm0.28$} & $\mathbf{80.94}${\scriptsize$\pm0.67$} & $\mathbf{73.13}${\scriptsize$\pm0.49$} & $\mathbf{58.87}${\scriptsize$\pm0.77$} & $\mathbf{77.24}${\scriptsize$\pm0.45$} & $\mathbf{13.52\%}$ \\
    \bottomrule
  \end{tabular}
\end{table}

\begin{figure}[t]
  \centering
  \includegraphics[width=0.50\linewidth]{figures/fig_accuracy_vs_severity.png}
  \caption{Top-1 accuracy against corruption severity, averaged over four
  corruption types and three seeds, with across-seed standard deviation as error
  bars. The curves do not cross at any severity: the ordering established on
  clean data is the ordering at every level of corruption.}
  \label{fig:severity}
\end{figure}

Absolute accuracy falls steeply with severity for every arm, from roughly $86$
to $89$\% clean down to $53$ to $59$\% at severity 5. The arms do not, however,
change places.

\subsection{Does the Clean Ranking Survive?}
\label{sec:ranking}

Table~\ref{tab:ranking} answers the title question directly.

\begin{table}[t]
  \centering
  \caption{Ranking of the four arms at each condition. No reordering occurs
  anywhere.}
  \label{tab:ranking}
  \small
  \begin{tabular}{lll}
    \toprule
    Condition & Ranking (best first) & Same as clean? \\
    \midrule
    Clean (severity 0) & CBAM $>$ SE $>$ none $>$ BAM & --- \\
    Severity 1 & CBAM $>$ SE $>$ none $>$ BAM & yes \\
    Severity 2 & CBAM $>$ SE $>$ none $>$ BAM & yes \\
    Severity 3 & CBAM $>$ SE $>$ none $>$ BAM & yes \\
    Severity 4 & CBAM $>$ SE $>$ none $>$ BAM & yes \\
    Severity 5 & CBAM $>$ SE $>$ none $>$ BAM & yes \\
    mCA (severities 1--5) & CBAM $>$ SE $>$ none $>$ BAM & yes \\
    Relative robustness drop & CBAM $>$ SE $>$ none $>$ BAM & yes \\
    \bottomrule
  \end{tabular}
\end{table}

\textbf{The clean-accuracy ranking survives.} It is preserved at all five
severities, on mean corruption accuracy, and on relative robustness drop.
Figure~\ref{fig:ranking} shows the same result as rank trajectories: a
reordering would appear as crossing lines, and none occurs. This is the null
outcome we pre-committed to reporting, and for this method set it means the
field's single clean-accuracy metric was not a misleading basis for choosing
among these blocks.

The relative-drop column is worth isolating, because it is the measure designed
to catch a spurious robustness advantage. CBAM loses $13.52\%$ of its clean
accuracy, SE $14.45\%$, the baseline $15.77\%$, and BAM $16.49\%$. The ordering
is unchanged, so CBAM's advantage under corruption is not an artefact of its
higher starting point: it degrades more slowly in proportional terms as well as
finishing higher.

This runs against our stated hypothesis. We expected the spatial gates of BAM
and CBAM to be more fragile than SE's global channel gate under blur and
contrast reduction. CBAM, which has a spatial gate, in fact degrades the most
slowly of all four arms. We report the hypothesis as refuted rather than
quietly restating it.

\begin{figure}[t]
  \centering
  \includegraphics[width=0.46\linewidth]{figures/fig_ranking_stability.png}
  \caption{Rank position of each arm at clean and at each corruption severity. A
  reordering would appear as crossing lines; none occurs, which is the paper's
  central negative result in its most direct form.}
  \label{fig:ranking}
\end{figure}

\subsection{What the Two Metrics Can and Cannot Resolve}
\label{sec:power}

A ranking is only meaningful if its adjacent pairs are separated by more than
seed noise, and here the two metrics differ in what they can resolve.

On clean accuracy the top two arms are $0.11$ pp apart, against a combined
across-seed standard deviation of $0.26$ pp. \textbf{Clean accuracy cannot
distinguish CBAM from SE at this budget.} On mean corruption accuracy the same
two arms are $0.92$ pp apart against a combined spread of $0.49$ pp, which does
exceed seed noise.

The corruption measurement therefore resolves a difference between the two best
arms that the clean measurement cannot. This is a distinct point from the
ranking result and cuts the other way: the corruption axis did not reorder
anything, but it discriminated more finely than the metric the field already
reports. A practitioner choosing between these two blocks on clean top-1 would
be choosing on a coin flip; the same choice on mCA is supported by the data.

The larger spreads in the table belong to the baseline and BAM arms
($\pm2.10$ and $\pm2.47$ pp clean, against $\pm0.18$ and $\pm0.08$ pp for SE and
CBAM). Attention appears to stabilise training across seeds at this budget as
well as improving the mean, though with three seeds we report this as an
observation rather than a claim.

\subsection{Per-Corruption Breakdown}

Averaging four corruption types can hide a type-specific skew, so
Table~\ref{tab:percorr} reports the breakdown.

\begin{table}[t]
  \centering
  \caption{Mean accuracy (\%) per corruption type, averaged over severities
  1--5 and three seeds.}
  \label{tab:percorr}
  \small
  \begin{tabular}{lcccc}
    \toprule
    Model & Brightness & Contrast & Defocus blur & Elastic transform \\
    \midrule
    none & $83.76$ & $59.66$ & $72.28$ & $74.38$ \\
    SE   & $86.83$ & $67.06$ & $74.63$ & $76.75$ \\
    BAM  & $78.66$ & $53.27$ & $69.40$ & $71.16$ \\
    CBAM & $\mathbf{87.33}$ & $\mathbf{67.53}$ & $\mathbf{76.05}$ & $\mathbf{78.05}$ \\
    \bottomrule
  \end{tabular}
\end{table}

The ordering is identical within every corruption type, so the aggregate is not
concealing a reversal. Contrast is by far the most damaging corruption, costing
the baseline $26$ pp relative to brightness, and it is also where the attention
arms separate most from the baseline: CBAM leads the baseline by $7.9$ pp on
contrast against $3.6$ pp on brightness. A global channel gate can rescale
channels whose activations have been uniformly compressed, which is a plausible
mechanism for that pattern, but our design does not isolate it and we offer it
as a hypothesis for future work rather than a finding.

\subsection{BAM Underperforms the Attention-Free Baseline}
\label{sec:bam}

BAM finishes last on every metric, $4.5$ pp below the baseline on clean accuracy
and $4.4$ pp below on mCA, with the widest across-seed spread of any arm, and the largest relative
degradation of the four (Table~\ref{tab:main}).

We are deliberately cautious about what this means. At a 10-epoch budget, an arm
that converges more slowly is indistinguishable from an arm that converges to a
worse optimum, and BAM is the heaviest module here, with dilated convolutions
and an extra batch-normalised MLP in every one of the eight blocks. BAM's mean
final-epoch training accuracy ($82.79$\%) is also the lowest of the four arms,
against $92.48$\% for CBAM and $87.90$\% for the baseline, which is consistent
with underfitting rather than with a robustness deficit. We therefore report
that \emph{BAM did not pay for itself at this budget}, and explicitly do not
claim that BAM is a worse module than the baseline in general. Section
\ref{sec:limitations} states what would settle it.

% ── 6. Limitations ────────────────────────────────────────────────────────────
\section{Limitations}
\label{sec:limitations}

\textbf{Training budget.} Ten epochs is matched across arms, which is what the
ranking claim requires, but it is short. Absolute accuracies would rise with a
longer schedule, and the ordering could in principle change --- most plausibly
for BAM, whose deficit (Section~\ref{sec:bam}) is consistent with underfitting.
The single most valuable extension is to rerun the identical grid at a longer
budget and check whether BAM's deficit closes.

\textbf{Three methods and one baseline.} The design supports per-severity
descriptive statements about whether the ranking held. It does not support a
general claim that attention \emph{type} predicts corruption robustness; four
points do not characterise a design space.

\textbf{Four of the benchmark's corruption types.} Averaging over brightness,
contrast, defocus blur and elastic transform could flatter or punish spatial
attention relative to the full benchmark. We report the per-type breakdown
alongside the mean so a skew is visible, and scope the finding to these four
types. All types in the source are ungated, so extension is inference-only and
requires no retraining.

\textbf{Insertion point.} Holding the slot fixed is what buys attribution to the
module, but it means no arm reproduces its own paper's placement. Absolute
accuracies are therefore not comparable to published results and none is
claimed.

\textbf{Three seeds, one dataset, one resolution.} Three seeds is enough to show
that the top-two clean gap is inside seed noise, but not to place tight
intervals on the arms with wider spread. Results are CIFAR-scale only.

% ── 7. Conclusion ─────────────────────────────────────────────────────────────
\section{Conclusion}
\label{sec:conclusion}

We trained three published attention blocks and an attention-free baseline into
one shared backbone, at one insertion point, under one budget and one seed set,
and re-ranked them under common image corruptions rather than on clean accuracy
alone.

The ranking survived. CBAM $>$ SE $>$ none $>$ BAM holds at every corruption
severity, on mean corruption accuracy, and on relative degradation. For this
method set, the clean-accuracy metric the literature reports was not a
misleading basis for choosing between these blocks, and our hypothesis that
spatial gates would prove more fragile than a global channel gate was refuted:
the arm with a spatial gate degraded most slowly.

Two findings qualify that. Clean accuracy could not resolve the top two arms
while mean corruption accuracy could, so the corruption axis was here the more
discriminative measurement even though it reordered nothing --- a reason to
report it that does not depend on it overturning anything. And BAM did not pay
for itself at this budget, finishing below the attention-free baseline, though
at ten epochs we cannot separate a worse optimum from slower convergence.

We also record that BAM's official implementation computes a multiplicative gate
fusion where its paper specifies an additive one. An implementation this widely
reused differing from its own paper on the fusion operator is worth the
field's attention independently of our results.

The cheapest extensions are clear: the remaining eleven corruption types are
ungated and require inference only, and rerunning the identical grid at a longer
budget would settle whether BAM's deficit is underfitting.

\bibliographystyle{plainnat}
\bibliography{refs}

\end{document}
